\section{Conclusion}\label{sec:conclusion}

We replicated a policy-gated multi-agent design for cloud pipeline governance with a live remote
language model, credible non-agent baselines, an adversarial evaluation against an independent oracle,
and closed-form sensitivity analyses, and we ran the same configuration under the deterministic mock
the original evaluation style relies on. The full agent system does beat static orchestration on
recovery time and cost, by margins smaller than and in one case directionally opposite to the
replicated paper's claims; it beats no non-agent comparator on speed; and its live evaluation and its
mock evaluation of the identical configuration disagree by roughly three orders of magnitude on
recovery time and by an equally large margin on decision quality. That last fact is the paper's
central claim: for this class of system, \emph{which} evaluation path produced a number is at least as
important as the number itself, and a claim carried from a mock evaluation to a live deployment
context is not a conservative extrapolation but a different experiment.

The system's architecture is Kirubakaran et al.'s, not ours. Our contribution is the measurement
discipline that makes claims about it falsifiable --- billing-accurate accounting, provenance that
refuses to silently merge incompatible arms, an oracle independent of the code under test, and seven
production defects found by holding the system to its own stated guarantees rather than to its test
suite --- and the results that discipline produced, reported as measured rather than reconciled to the
system's advantage. We take the practical implication to be narrow and specific rather than a verdict
on LLM agents in operations generally: before a governance claim about a policy-gated agent system is
acted on, ask what evaluated it, and treat a mock number and a live number as two different findings
until both exist.
